\clearpage\thispagestyle{empty}
\pdfbookmark[0]{Quellen zu Gauß’ Bericht über den neuen Planeten vom Oktober 1804}{nachrichten1804-sources}
\begin{center}\textbf{Quellen zu Gauß’ Bericht über den neuen Planeten vom Oktober 1804}\end{center}
\noindent Diese KI-Übertragung in heutiges Deutsch folgt Carl Friedrich Gauß’ Bericht in den \href{https://gdz.sub.uni-goettingen.de/id/PPN236007467}{\textit{Werken}, Band~VI} (Göttingen, 1874), S.~253--256. Das \href{https://digipress.digitale-sammlungen.de/view/bsb10531607_00329_u001}{\textit{Braunschweigische Magazin}}, Stück~40 vom 6. Oktober 1804, nennt Gauß als Verfasser in \href{https://www.digitale-sammlungen.de/en/view/bsb10531607?page=329}{Sp.~621}. Der Erstdruck endet mit einer Ephemeride in \href{https://www.digitale-sammlungen.de/en/view/bsb10531607?page=332}{Sp.~628}. Die Werkausgabe verweist auf die \hyperlink{native-258}{Ephemeride in Franz Xaver von Zachs Bericht auf S.~253}.
\par\smallskip\noindent Karl Ludwig Harding entdeckte Juno. Der von Harding verwendete Sternkatalog war Jérôme de Lalandes \href{https://archive.org/download/histoireclestef00lalagoog/page/n12.jpg}{\textit{Histoire céleste française}} (Paris, 1801). Gauß nennt Michel Lefrançais Lalande für seine Beobachtungen und John Flamsteed für dessen Sternverzeichnis.
\par\smallskip\noindent Übertragung und Satz wurden mit KI erstellt. Die Digitalisate der Zeitschrift stammen von der Bayerischen Staatsbibliothek.
\clearpage
